\section*{Overview}

2-SAT is the point where Boolean satisfiability stops feeling like brute force and starts feeling like graph theory.
Once every clause has exactly two literals, the whole problem becomes an implication graph plus strongly connected
components.

\section*{Problem-Driven Motivation}

Many contest statements look like this after modeling:

\begin{itemize}
  \item each object has exactly two possible states,
  \item some pairs of chosen states are incompatible,
  \item some pair of states must contain at least one true choice.
\end{itemize}

The naive space is \(2^n\), which is hopeless once \(n\) reaches \(10^5\). But the constraints are not arbitrary SAT
clauses. They are pairwise logical constraints, and that extra structure is enough to solve them in linear time.

\section*{Recognition Pattern}

2-SAT is the right direction when:

\begin{itemize}
  \item every decision is binary,
  \item constraints can be written as clauses \((a \lor b)\),
  \item the problem contains implications like ``if this state is chosen, that state is forbidden / forced,''
  \item the intended model is a large compatibility system rather than a numeric optimization.
\end{itemize}

Classic signals are:

\begin{itemize}
  \item choose one of two colors / directions / schedules,
  \item at least one of two options must be taken,
  \item two specific literal choices cannot both happen.
\end{itemize}

\section*{Derivation}

The key equivalence is:
\[
(a \lor b) \iff (\lnot a \Rightarrow b)\ \text{and}\ (\lnot b \Rightarrow a).
\]

So instead of treating clauses directly, build a directed graph on literals:

\begin{itemize}
  \item one node for each variable being true,
  \item one node for each variable being false.
\end{itemize}

Then every clause adds two implications.

Why do SCCs solve the problem?

\begin{itemize}
  \item If \(x\) and \(\lnot x\) lie in the same SCC, then each implies the other, so the formula forces \(x\) to be
  simultaneously true and false. That is impossible.
  \item Otherwise, the SCC condensation graph is a DAG. Processing SCCs in reverse topological order lets us choose a
  literal before its negation whenever possible, producing a valid assignment.
\end{itemize}

\section*{Worked Problem}

\paragraph{Problem.}
Each worker must be assigned to one of two shifts: day or night. Some pairs of concrete choices are incompatible. For
example:

\begin{itemize}
  \item if worker \(u\) takes day, worker \(v\) cannot take night,
  \item at least one of worker \(a\) taking day or worker \(b\) taking day must happen.
\end{itemize}

\paragraph{Why naive fails.}
Trying all shift assignments is \(2^n\).

\paragraph{Modeling.}
Let variable \(x_i\) mean ``worker \(i\) takes day.'' Then:

\begin{itemize}
  \item ``worker \(u\) day and worker \(v\) night cannot both happen'' becomes
  \[
  (\lnot(x_u \land \lnot x_v)) \iff (\lnot x_u \lor x_v),
  \]
  \item ``at least one of \(x_a\) or \(x_b\)'' is already a 2-SAT clause:
  \[
  (x_a \lor x_b).
  \]
\end{itemize}

\paragraph{Graph construction.}
Every clause \((p \lor q)\) contributes:
\[
\lnot p \Rightarrow q,\qquad
\lnot q \Rightarrow p.
\]

\paragraph{Final algorithm.}

\begin{itemize}
  \item build the implication graph,
  \item compute SCCs,
  \item reject if any variable and its negation share an SCC,
  \item otherwise assign values by SCC order.
\end{itemize}

\section*{Implementation Reasoning}

The only implementation detail that really matters is a clean literal encoding.

The code uses:

\begin{itemize}
  \item node \(2x\) for ``variable \(x\) is true,''
  \item node \(2x+1\) for ``variable \(x\) is false,''
  \item negation by \(\texttt{literal} \oplus 1\).
\end{itemize}

That makes the clause and implication helpers short and hard to misuse.

The solver exposes:

\begin{itemize}
  \item \texttt{add\_or},
  \item \texttt{add\_implication},
  \item \texttt{force\_true},
  \item \texttt{satisfiable()}.
\end{itemize}

\section*{Correctness Intuition}

The implication graph captures the exact situations where one false literal forces another literal to become true. SCCs
detect contradictions because mutual reachability means mutual forcing. Once contradictions are gone, assigning literals
in reverse SCC topological order guarantees that choosing a literal never violates an implication to an already-fixed
literal.

\section*{Complexity Analysis}

For \(n\) variables and \(m\) clauses:

\begin{itemize}
  \item vertices: \(2n\),
  \item edges: \(2m\) plus any extra forced implications,
  \item total time: \(O(n+m)\),
  \item memory: \(O(n+m)\).
\end{itemize}

\section*{Common Pitfalls}

\begin{itemize}
  \item Mixing variable indices with literal indices.
  \item Extracting the assignment from DFS order instead of SCC order.
  \item Forgetting that ``not both'' is a clause, not a direct edge, until translated correctly.
  \item Trying to force constraints that are not reducible to 2-literal clauses.
\end{itemize}

\section*{Variants and Failure Modes}

\begin{itemize}
  \item ``At most one'' over many choices needs extra encoding; it is not one 2-SAT clause by itself.
  \item General SAT does not collapse to SCCs.
  \item Sometimes a graph-coloring or bipartite formulation is simpler than 2-SAT, even if 2-SAT is possible.
\end{itemize}

\section*{Practice Problems}

\begin{itemize}
  \item Binary scheduling with pairwise incompatibilities.
  \item Orientation or color choices with logical constraints.
  \item Giant Pizza and similar implication-graph problems.
\end{itemize}

\section*{References}

\begin{itemize}
  \item \href{https://cp-algorithms.com/graph/2SAT.html}{cp-algorithms: 2-SAT}.
  \item \href{https://usaco.guide/adv/SCC}{USACO Guide: SCC and 2-SAT}.
  \item \href{https://codeforces.com/catalog}{Codeforces Catalog}.
\end{itemize}
